\chapter{DOCUMENT METADATA VÀ REVISION HISTORY}

\begin{ChapterObjectives}
\item Xác định các metadata quản trị tài liệu cần có cho bộ technical manual của dự án.
\item Ghi nhận version, owner, reviewer, approval status, và change log theo cách phù hợp với trạng thái hiện tại của project.
\item Tạo một baseline revision chapter để các lần cập nhật sau có thể nối tiếp thay vì quản lý tài liệu rời rạc.
\item Giúp tài liệu gần hơn với một bản release-ready handover manual thay vì chỉ là một tập hợp chapter kỹ thuật.
\end{ChapterObjectives}

\section{Vai trò của document metadata}

Khi technical manual đã bắt đầu có nhiều chapter, việc bổ sung metadata quản trị tài liệu là cần thiết để trả lời các câu hỏi sau:

\begin{enumerate}[label=-]
\item tài liệu này thuộc project nào;
\item phiên bản tài liệu hiện tại là gì;
\item ai đang là người sở hữu hoặc biên soạn chính;
\item ai sẽ review;
\item trạng thái phê duyệt hiện tại ra sao;
\item những thay đổi lớn nào đã được ghi nhận qua từng revision.
\end{enumerate}

Các thông tin này không thay thế commit history, nhưng chúng giúp người đọc PDF có một lớp ngữ cảnh quản trị tài liệu ngay trong bản phát hành. Điều này đặc biệt hữu ích khi tài liệu được in ra, gửi offline, hoặc dùng như hồ sơ bàn giao kỹ thuật.

\section{Document metadata baseline}

Bảng \ref{tab:document_metadata} là baseline metadata cho bộ technical manual hiện tại. Các trường đã chắc chắn thì được điền trực tiếp; các trường chưa được chốt chính thức thì được giữ ở trạng thái reviewable để cập nhật sau.

\begin{table}[H]
	\centering
	\renewcommand{\arraystretch}{1.3}
	\caption{Document metadata baseline}
	\label{tab:document_metadata}
	\begin{tabularx}{\linewidth}{|>{\RaggedRight\arraybackslash}p{3.5cm}|>{\RaggedRight\arraybackslash}X|}
		\hline
		\textbf{Metadata Field} & \textbf{Current Value} \\ \hline
		Project name & Energy Consumption Reporting System \\ \hline
		Document type & Project technical manual / engineering reference document \\ \hline
		Document scope & Architecture, data sources, repository/service mapping, business rules, KPI rules, output structure, release workflow, deployment and operations, scheduling, troubleshooting, glossary, and naming conventions \\ \hline
		Current document version & \texttt{v1.0} \\ \hline
		Document status & Internal release baseline \\ \hline
		Primary owner / author & Tuyen Tran \\ \hline
		Reviewer & Reviewer assignment pending \\ \hline
		Approval status & Pre-handover sign-off pending \\ \hline
		Canonical LaTeX source root & \TablePath{docs/latex/templates/project_technical_documentation_template/} \\ \hline
		Canonical PDF output & \TablePath{docs/output/project_technical_documentation_latest.pdf} \\ \hline
		Expected dated PDF output & \TablePath{docs/output/project_technical_documentation_YYYYMMDD.pdf} \\ \hline
		Revision control source & Git commits + chapter-level revision history table in this document \\ \hline
	\end{tabularx}
\end{table}

\section{Versioning convention}

Phiên bản tài liệu nên được quản lý độc lập nhưng có liên hệ rõ với trạng thái implementation của project. Từ mốc release baseline hiện tại, một convention đơn giản nhưng an toàn là:

\begin{enumerate}[label=-]
\item dùng \texttt{v1.0} cho bản technical manual release baseline đầu tiên;
\item tăng patch hoặc minor version khi có thêm chapter, rule, hoặc thay đổi đáng kể ở mức tài liệu;
\item chỉ dùng nhãn draft/review trong revision history khi tài liệu quay lại một vòng chỉnh sửa chưa sẵn sàng phát hành;
\item nếu cần, gắn thêm commit SHA hoặc release candidate note trong revision history thay vì nhồi quá nhiều ý nghĩa vào tên file PDF.
\end{enumerate}

Mục tiêu của convention này là giúp người đọc biết tài liệu đang ở mức draft, review, hay release baseline, mà không làm lẫn version của tài liệu với version runtime của application.

\section{Author, reviewer, và approval responsibilities}

Tài liệu kỹ thuật chất lượng tốt cần có vai trò rõ ràng giữa người viết, người rà soát, và người phê duyệt. Bảng \ref{tab:document_roles} tóm tắt responsibilities khuyến nghị.

\begin{table}[htbp]
	\centering
	\renewcommand{\arraystretch}{1.3}
	\caption{Document roles and responsibilities}
	\label{tab:document_roles}
	\begin{tabularx}{\linewidth}{|>{\RaggedRight\arraybackslash}p{2.7cm}|>{\RaggedRight\arraybackslash}X|}
		\hline
		\textbf{Role} & \textbf{Responsibility} \\ \hline
		Primary owner / author & Cập nhật nội dung tài liệu, giữ tài liệu bám implementation, và đảm bảo các chapter quan trọng luôn có trạng thái hợp lý. \\ \hline
		Reviewer & Rà soát độ đúng kỹ thuật, tính nhất quán giữa tài liệu và runtime, kiểm tra business rules, KPI rules, và output conventions. \\ \hline
		Approver / sign-off owner & Xác nhận tài liệu đủ tốt để dùng cho review, handover, hoặc phát hành nội bộ; quyết định khi nào đổi trạng thái tài liệu từ draft sang approved baseline. \\ \hline
		Automation / OpenClaw support & Hỗ trợ build PDF, kiểm tra cross-reference, phát hiện drift giữa docs và implementation, và giữ workflow Documentation-as-Code hoạt động ổn định. \\ \hline
	\end{tabularx}
\end{table}

Ở trạng thái hiện tại, owner đã rõ. Reviewer và approver vẫn chưa được gán tên chính thức, nhưng trạng thái pending hiện đã được diễn đạt rõ ràng để tài liệu có thể dùng như một bản internal release baseline mà không gây hiểu nhầm là đã sign-off xong.

\section{Approval and sign-off flow}

Một luồng phê duyệt tài liệu kỹ thuật tối thiểu có thể là:

\begin{enumerate}[label=-]
\item owner cập nhật source LaTeX/Markdown liên quan;
\item build lại PDF kỹ thuật và rà soát formatting / cross-reference;
\item reviewer kiểm tra độ đúng của nội dung kỹ thuật;
\item nếu tài liệu được dùng cho milestone hoặc handover, approver xác nhận trạng thái đủ tốt để phát hành nội bộ;
\item revision history được cập nhật cùng với kết quả review hoặc sign-off note.
\end{enumerate}

Luồng này không cần quá nặng ở giai đoạn draft, nhưng nên tồn tại rõ ràng để khi tài liệu trưởng thành hơn, nhóm không phải tự phát minh lại quy trình quản trị tài liệu từ đầu.

\section{Revision history baseline}

Bảng \ref{tab:revision_history} là baseline revision history của bộ technical manual hiện tại.

\begin{table}[htbp]
	\centering
	\renewcommand{\arraystretch}{1.3}
	\caption{Technical manual revision history baseline}
	\label{tab:revision_history}
	\begin{tabularx}{\linewidth}{|>{\RaggedRight\arraybackslash}p{2.2cm}|>{\RaggedRight\arraybackslash}p{2.0cm}|>{\RaggedRight\arraybackslash}p{2.6cm}|>{\RaggedRight\arraybackslash}X|}
		\hline
		\textbf{Version} & \textbf{Date} & \textbf{Status} & \textbf{Change Summary} \\ \hline
		\texttt{v1.0} & 2026-05-12 & Internal release baseline & Chốt một release baseline duy nhất cho technical manual hiện tại, gom các draft trước đó thành mốc phát hành \texttt{v1.0} để PDF handover/review dễ đọc và không còn bảng lịch sử draft quá dài. \\ \hline
	\end{tabularx}
\end{table}

Mục đích của bảng này không phải thay thế Git, mà là giúp người đọc PDF hiểu mốc trưởng thành của tài liệu mà không cần truy cập repository ngay lập tức.

\section{Change log policy}

Change log của tài liệu nên tuân theo các nguyên tắc sau:

\begin{enumerate}[label=-]
\item chỉ ghi những thay đổi có ý nghĩa ở mức tài liệu hoặc workflow, không cần ghi mọi edit nhỏ;
\item khi một chapter mới được thêm hoặc một rule cốt lõi bị đổi, revision history phải được cập nhật;
\item nếu technical PDF được dùng cho handover hoặc internal release review, nên ghi rõ version đang được xem xét;
\item nếu runtime thay đổi làm một chapter cũ không còn đúng, phải ưu tiên cập nhật chapter đó trước khi coi PDF là usable reference.
\end{enumerate}

\section{Document readiness checklist}

Trước khi xem technical PDF là đủ trưởng thành cho một vòng review nghiêm túc, nên rà soát tối thiểu các điểm sau:

\begin{enumerate}[label=-]
\item document metadata đã phản ánh đúng project name, source root, output location, và document status;
\item version hiện tại đã rõ là draft, review, hay release baseline;
\item revision history đã phản ánh mốc release baseline hiện tại;
\item reviewer và approval path đã được xác định hoặc ít nhất được để ở trạng thái pending assignment có chủ đích;
\item tài liệu đã build pass, cross-reference sạch, và PDF output đã được cập nhật theo version mới nhất.
\end{enumerate}

\section{Tổng kết chương}

Chapter này bổ sung lớp quản trị tài liệu cho technical manual: document metadata, versioning convention, role mapping, approval flow, revision history, và change log policy. Sau chapter này, bộ tài liệu không chỉ đầy đủ về mặt kỹ thuật mà còn tiến thêm một bước về mặt document governance và readiness cho review hoặc handover.
